\documentclass[t]{beamer}
\usepackage{pgfpages}
\usepackage[ngerman]{babel}
\usepackage{textcomp}
\usepackage{subfigure}
\usepackage{graphicx}
\usepackage{fancyhdr}
\usepackage{tgpagella}
\usepackage{textpos}
\usepackage{hyperref}
\usepackage{tcolorbox}
\usepackage{tikz}
\usetikzlibrary{positioning}
\usepackage{pdfpcnotes}

\usepackage{fontspec}
\setmainfont{Titillium Web}
\setsansfont{Titillium Web}

\usetheme{Dresden}
\useoutertheme{infolines}
\usecolortheme{fsfw}

\setbeamerfont*{quote}{family=\sffamily}
\addtobeamertemplate{quote begin}{}{\begin{tcolorbox}[colback=fsfw-gray,colframe=fsfw-violet!20!fsfw-gray]}
\addtobeamertemplate{quote end}{\end{tcolorbox}}{}

\addtobeamertemplate{block example begin}{\begin{tcolorbox}[colback=fsfw-gray,colframe=fsfw-lemon!20!fsfw-gray]}{}
\addtobeamertemplate{block example end}{}{\end{tcolorbox}}

\title{Linux Presentation Day April 2016}
\author{\texttt{https://fsfw-dresden.de/lpd}}
\date[30.04.2016]{30. April 2016}

\newcommand{\cmd}[1]{\texttt{#1}}

% TODO: fsfw Farbschema
% LPD Logo in die Ecke
\addtobeamertemplate{frametitle}{}{%
    \begin{textblock*}{130mm}(.82\textwidth,-0.5cm)
      \includegraphics[width=1.8cm]{pdf-gen/lpd-logo-small.pdf}
\end{textblock*}
\vspace*{1em}}

\hyphenation{
  Regel
}

\AtBeginSection[] % Do nothing for \section*
{
  \begin{frame}<beamer>
    \frametitle{Gliederung}
    \tableofcontents[currentsection]
  \end{frame}
}


%% \setbeameroption{show notes on second screen=right}
\setbeamertemplate{navigation symbols}{}%remove navigation symbols
\setbeamertemplate{note page}{%
  \insertnote%
}

\AtBeginNote{
    \let\enumerate\itemize%
    \let\endenumerate\enditemize%
    \hspace{-0.5cm}\begin{minipage}{\textwidth}
        %% \begin{columns}
        %%     \column{0.4\textwidth}
        %%     \scriptsize
        %%     Current slide: \\
        %%     \insertslideintonotes{0.4}
        %%     \column{0.59\textwidth}
}
\AtEndNote{
        %% \end{columns}
    \end{minipage}
}
\newcommand{\noteitem}[1]{\pnote{#1}}

\begin{document}
\begin{frame}
  \begin{center}%
% LPD Logo
\includegraphics[width=0.8\textwidth]{pdf-gen/lpd-logo.pdf}\\%
\structure{\Large Linux Presentation Day 2016}\vspace{0.5cm}\\
{\includegraphics[height=1cm]{pdf-gen/dkn_logo.pdf}\hspace{0.5cm}\includegraphics[height=1cm]{pdf-gen/c3d2-logo.pdf}\\\includegraphics[height=1cm]{pdf-gen/tux.pdf}\hspace{0.5cm}\includegraphics[height=1cm]{pdf-gen/fsfw-logo.pdf}}
  \end{center}
\end{frame}

\begin{frame}
  \frametitle{Gliederung}
  \tableofcontents
\end{frame}

\section{Einführung und Vorstellung}
\begin{frame}
  \frametitle{Wer sind wir?}
  \begin{itemize}
  \item FSFW – Hochschulgruppe für Freies Wissen und Freie Software
  \item LUG DD – Linux-User-Group Dresden
  \item c3d2 – Chaos Computer Club Dresden
  \item Datenkollektiv
  \item StuRa der HTW (Räume und Budget)
  \item IT4R – IT for refugees
  \end{itemize}
  % Pro Initiative eine Slide (und eine Übersichtsslide)
  % URLs
\end{frame}

\begin{frame}
  \frametitle{Freie Software Freies Wissen}
  \begin{columns}
    \begin{column}[T]{0.45\textwidth}
      \includegraphics[width=\textwidth]{pdf-gen/fsfw-logo-with-text.pdf}
    \end{column}
    \begin{column}{0.45\textwidth}
      \begin{itemize}
      \item Hochschulgruppe an der TU Dresden
      \item Mitglieder von der HTW und Mitglieder ohne Hochschulaffiliation
      \end{itemize}
    \end{column}
  \end{columns}
\end{frame}

\begin{frame}
  \frametitle{Linux User Group Dresden}
  \begin{columns}
    \begin{column}[T]{0.45\textwidth}
      \includegraphics[width=\textwidth]{pdf-gen/tux.pdf}
    \end{column}
    \begin{column}{0.45\textwidth}
      \begin{itemize}
      \item Loser Verband von Linuxinteressierten
      \item Zweiwöchentliche Treffen
      \item Bietet gerne Hilfe für Neueinsteiger
      \end{itemize}
    \end{column}
  \end{columns}
\end{frame}

\begin{frame}
  \frametitle{Chaos Computer Club Dresden}
  \begin{center}
    \includegraphics[width=0.5\textwidth]{pdf-gen/c3d2-logo.pdf}
  \end{center}
  \begin{itemize}
  \item Örtlicher Ableger des CCC e.V.
  \item Hackerspace
  \item Interesse an Datenschutz und Datensicherheit
  \end{itemize}
\end{frame}

\begin{frame}
  \frametitle{Datenkollektiv}
  \includegraphics[width=\textwidth]{pdf-gen/dkn_logo.pdf}
  \begin{itemize}
  \item Betrieb von Mailservern seit Debatte um Vorratsdatenspeicherung
  \item Heute auch mehr Dienste
  \item GbR seit 2014
  \item Fokus auf lokale Verankerung, Hardwarerecycling und selbstbestimmten Umgang mit Technik
  \item „semikommerziell“
  \end{itemize}
\end{frame}

\begin{frame}
  \frametitle{Studentenrat der HTW Dresden}
  \begin{itemize}
  \item Hochschulweite studentische Vertretung
  \item Ressort Datenkultur
  \end{itemize}
\end{frame}

\begin{frame}
  \frametitle{IT 4 Refugees}
  \begin{itemize}
  \item Schulung von Aktivisten (Datensicherheit,
    E-Mail-Verschlüsselung)
  \item Bereitstellen von Internetverbindungen für
    Flüchtlingsunterkünfte
  \end{itemize}
\end{frame}

\begin{frame}
  \frametitle{Warum machen wir das?}
  \begin{itemize}
  \item weil heute der internationale Linux-Presentation-Day ist
  \item aktuelles Interesse an Linux-Systemen (Windows 10, Steam
    für Linux, Edward Snowden)
  \item wir verwenden es selbst gerne
  \end{itemize}
\end{frame}

\section{Über Linux und Freie Software}
\begin{frame}
  \noteitem{Quote vorlesen}
  \noteitem{folgende Struktur erklären (Grundlegende Begriffe OS, FLOSS, dann zurück zu Linux)}
  \frametitle{Was ist Linux?}
  \begin{quote}
    Als Linux oder GNU/Linux […] bezeichnet man in
    der Regel freie, unix-ähnliche Mehrbenutzer-Betriebssysteme, die
    auf dem Linux-Kernel und wesentlich auf GNU-Software
    basieren.\footnote{\url{https://de.wikipedia.org/wiki/Linux}}
  \end{quote}
  \begin{center}
    \Large … das bedarf Erläuterung.
  \end{center}
\end{frame}

\begin{frame}
  \noteitem{Bloß nicht an Scheduling, ... aufhängen!}
  \noteitem{Userland genauer ausführen}
  \frametitle{Was ist ein Betriebssystem?}
  \begin{columns}
    \begin{column}[c]{0.49\textwidth}
      \centering
      \begin{tikzpicture}[node distance=-1pt]
        \tikzset{
          part/.style={
            draw=fsfw-violet!20!fsfw-gray,
            fill=fsfw-gray,
            rectangle,
            line width=1pt,
            inner sep=5pt,
          }
        }
        \node[part,minimum width=4em] (hardware) {Gerät};
        \node[part,below=of hardware,minimum width=5em] (kern) {Kernel};
        \node[part,below=of kern,minimum width=6em] (svc) {Dienste};
        \node[part,below=of svc,minimum width=7em] (userland) {Userland};
      \end{tikzpicture}
    \end{column}
    \begin{column}[c]{0.5\textwidth}
      \begin{block}{Kernel}
        macht die Drecksarbeit – stellt die Umgebung für Programme bereit
        \begin{itemize}
        \item Scheduling, Speicherverwaltung, Hardwaretreiber, Interrupts,
          DMA, Netzwerk, …
        \end{itemize}
      \end{block}
      \begin{block}{Userland}
        \begin{itemize}
        \item Graphische Oberfläche
        \item Standardwerkzeuge
        \item Standardbibliotheken
        \end{itemize}
      \end{block}
    \end{column}
  \end{columns}
\end{frame}

\begin{frame}
  \noteitem{Lizenzdialog!}
  \noteitem{Entwicklungsmodell}
  \noteitem{FSF und GNU}
  \frametitle{Was bedeutet „Freie Software“?}
  \begin{columns}
    \begin{column}{0.45\textwidth}
      \begin{itemize}
      \item Bewegung mitbegründet von Richard M. Stallman
      \end{itemize}
      \begin{block}{Grundfreiheiten für Software}
        \begin{enumerate}
        \item Zu jedem Zwecke ausführen
        \item Untersuchen und Verändern
        \item Weitergeben
        \item Weitergeben von Änderungen
        \end{enumerate}
      \end{block}
    \end{column}
    \begin{column}[T]{0.45\textwidth}
      %% \includegraphics[width=\textwidth]{img-src/vlc-license.png}\\
      %% \includegraphics[width=\textwidth]{pdf-gen/gplv3-logo.pdf}
      \includegraphics[width=\textwidth]{img-src/rms.jpg}
    \end{column}
  \end{columns}
  % Entwicklungsmodell – OpenSource
\end{frame}

\begin{frame}
  \frametitle{Was ist Linux? – for real}
  \begin{columns}
    \begin{column}[T]{0.45\textwidth}
      \includegraphics[width=\textwidth]{img-src/linux-gnu.png}
    \end{column}
    \begin{column}{0.45\textwidth}
      \begin{itemize}
      \item bezeichnet eigentlich den Linux-Kernel
      \item im allgemeinen Sprachgebrauch: ganzes Betriebssystem
      \item Freie Software – größtenteils
      \item große Vielfalt an Ausprägungen (Server, Telefone, diverse
        Desktop-Systeme)
      \end{itemize}
    \end{column}
  \end{columns}
\end{frame}

\begin{frame}
  \frametitle{Geschichte von Linux}
  \begin{columns}
    \begin{column}[T]{0.4\textwidth}
      \includegraphics[width=\textwidth]{img-src/linus.jpg}
    \end{column}
    \begin{column}{0.5\textwidth}
      \begin{itemize}
      \item Entwickelt von Linus Torvalds
      \item Hobbyprojekt
      \item Kein vergleichbares Projekt zu der Zeit
      \item Großes Interesse unter Entwicklern
      \end{itemize}
    \end{column}
  \end{columns}
\end{frame}

% Geschichte von Linux Slide

\section{Begriffe}
\begin{frame}
  \noteitem{Software umfasst für gewöhnlich allgemeine Freie Software, Freie Software die für die Distribution gefertigt wurde und proprietäre Software.}
  \frametitle{Distribution}
  \includegraphics[angle=-90,width=\textwidth]{pdf-gen/distros.pdf}
  \begin{itemize}
  \item Sammlung von fertig paketierter Software und Konfiguration.
  \item Ansprechpartner bei Problemen
  \item Stellt (Sicherheits-)Updates bereit
  \end{itemize}
  \begin{exampleblock}{Beispiele}
    Fedora, Debian, Linux Mint, Ubuntu, ArchLinux
  \end{exampleblock}
\end{frame}

\begin{frame}
  \noteitem{Graphische Oberfläche: Fensterverwaltung, Menü zum Starten von Anwendungen}
  \frametitle{Desktop-Umgebung}
  \begin{columns}
    \begin{column}[T]{0.45\textwidth}
      \includegraphics[width=\textwidth]{img-src/cinnamon.png}
    \end{column}
    \begin{column}{0.45\textwidth}
      \begin{itemize}
      \item Graphische Oberfläche
      \item Standardanwendungen (PDF-Anzeigetool, Taschenrechner, Mine Sweeper, Browser, E-Mail, …)
      \end{itemize}
      \begin{exampleblock}{Beispiele}
        GNOME, Cinnamon, KDE, Ubuntu Unity, …
      \end{exampleblock}
    \end{column}
  \end{columns}
\end{frame}

\begin{frame}
  \frametitle{Datei-Manager}
  \begin{columns}
    \begin{column}[T]{0.45\textwidth}
      \includegraphics[width=\textwidth]{img-src/file-manager.png}
    \end{column}
    \begin{column}{0.45\textwidth}
      \begin{itemize}
      \item Graphisches Werkzeug zum Finden und Verwalten von Dateien
      \item vgl. Windows Explorer, OS X Finder
      \end{itemize}
    \end{column}
  \end{columns}
\end{frame}

\section{Gründe für und gegen Linux}
\begin{frame}
  \frametitle{Warum Freie Software?}
  \begin{itemize}
  \item Vier Freiheiten
    % wichtige anmerkung: => Verifizierbarkeit
  \item keine Lizenzgebühren
  \item Unterstützung älterer Hardware
  \item ethische Argumente
  \item Zugänglichkeit für Bildung
  \item Nachhaltigkeit
  \item Sicherheit${}^\text{[citation needed]}$
  \end{itemize}
\end{frame}

\begin{frame}
  \frametitle{Kontrolle über Updates}
  \begin{itemize}
  \item Updates werden installiert, wenn und wann der Nutzer das möchte
  \item Updates brauchen in der Regel keinen Neustart
  \item Sicherheitsupdates können gesondert installiert werden
  \item kein verlängertes Herunterfahren durch Updates
  \end{itemize}
\end{frame}

\begin{frame}
  \frametitle{Konfigurierbarkeit und Anpassbarkeit}
  \begin{itemize}
  \item große Auswahl aus heterogenem Ökosystem
  \item für jeden etwas dabei (alte, schwächere Hardware; minimale
    Nerd-Systeme; ...)
  \item die meiste Software ist sehr gut konfigurierbar
  \end{itemize}
\end{frame}

\begin{frame}
  \frametitle{Privacy}
  \begin{columns}
    \begin{column}[T]{0.45\textwidth}
      \includegraphics[width=\textwidth]{pdf-gen/privacy.pdf}
    \end{column}
    \begin{column}{0.45\textwidth}
      \begin{itemize}
      \item telefoniert nicht nach Hause
      \item Festplattenverschlüsselung ist typischerweise einfach (ein
        Klick bei der Installation und ein Passwort angeben)
      \end{itemize}
    \end{column}
  \end{columns}
\end{frame}

\begin{frame}
  \noteitem{Sicherheit für Linux zum Glück kein Neuland}
  \frametitle{Sicherheit}
  \begin{columns}
    \begin{column}[T]{0.45\textwidth}
      \includegraphics[width=\textwidth]{img-src/neuland.jpg}
    \end{column}
    \begin{column}{0.45\textwidth}
      \begin{itemize}
      \item Vorteil durch kleineren Marktanteil im Desktopbereich
      \item Vorteil durch großen Marktanteil auf Serversystemen
      \item Community reagiert typischerweise schnell auf
        Sicherheitslücken
        % mündlich: ob man auch Updates bekommt ...
      \end{itemize}
    \end{column}
  \end{columns}
\end{frame}

\begin{frame}
  \frametitle{Paketverwaltung}
  \begin{columns}
    \begin{column}[T]{0.45\textwidth}
      \includegraphics[width=\textwidth]{img-src/pkg-manager.png}
    \end{column}
    \begin{column}{0.45\textwidth}
      \begin{itemize}
      \item Programme komfortabel und sicher installieren
      \item von der Distribution bereitgestellt
      \item für viele \emph{das} Killerfeature
      \item App-Store auf Steroiden – seit 1998
        % mündlich: Abhängigkeiten, Deduplizierung, Katzenbilder
        % Programme für $alles direkt verfügbar ohne windige
        % Download-Websites
      \end{itemize}
    \end{column}
  \end{columns}
\end{frame}

\begin{frame}
  \frametitle{Nutzerfreundlichkeit}
  \begin{itemize}
  \item schlechter Ruf nicht mehr gerechtfertigt
  \item mit modernen Desktopumgebungen mindestens so komfortabel wie
    Windows
  \item keine Angst vor der Konsole
    \begin{itemize}
    \item Kenntnis nicht zwingend notwendig
    \item macht das Leben für erfahrene Benutzer aber viel einfacher
    \end{itemize}
  % mündlich: it's a feature, not a bug
  % Symbolbolbild: Terminal
  % Eventuell eigene Slide für die Konsole
  \end{itemize}
\end{frame}

\begin{frame}
  \frametitle{Probleme mit Linux}
  \begin{itemize}
  \item exotische Hardware (NFC, seltener: WLAN, sehr neue Graphikkarten)
  \item Graphikleistung oft schlechter ($\rightarrow$ Spiele)
  \item Verfügbarkeit von spezieller Software (z.\,B. CAD)
  \item Accessibility
    % mündlich: NFC erläutern
  \end{itemize}
\end{frame}

\begin{frame}
  \noteitem{einziges praktisches Problem: Windows zerlegt den Bootloader bei größeren Updates (das muss man dann manuell beheben)}
  \frametitle{Linux neben Windows: Dual-Boot}
  \begin{columns}
    \begin{column}[T]{0.45\textwidth}
      \includegraphics[width=\textwidth]{img-src/grub.png}
    \end{column}
    \begin{column}{0.45\textwidth}
      \begin{definition}{Dual-Boot}
        Zwei (oder mehr) Betriebssysteme auf einem Rechner installiert haben. Beim Einschalten wählt man, welches gestartet werden soll.
      \end{definition}
    \end{column}
  \end{columns}
  \begin{itemize}
  \item im Prinzip kein Problem
  \item man bekommt dann ein Bootmenü zum Auswählen
  \end{itemize}
  % TODO: Screenshot von grubd
  % zwei Folien: Dual Boot/Alternativen zu Dual Boot
  % zweite Slide: VM und wine erklären
  % Programmnamen Mark-Upen
\end{frame}

\begin{frame}
  \noteitem{pro/contra wine/vm (wine => performance, vm => stabilität)}
  \frametitle{Einzelne Windows-Anwendungen verwenden}
  \begin{columns}
    \begin{column}[T]{0.35\textwidth}
      \includegraphics[width=\textwidth]{pdf-gen/wine-logo.pdf}
    \end{column}
    \begin{column}{0.55\textwidth}
      \begin{itemize}
      \item viele Windows-Programme laufen auch unter \cmd{wine} (ältere Spiele, Bürosoftware, ...)
      \item falls (Graphik-)Leistung oder direkter Hardwarezugriff nicht wichtig sind: Windows in einer virtuellen Maschine
      \end{itemize}
    \end{column}
  \end{columns}
\end{frame}

\section{Programm}
\begin{frame}
  \frametitle{Was gibt es jetzt noch?}
  \begin{itemize}
  \item Linux-Install-Party
  \item Computer mit Linux zum Ausprobieren
  \item Vorstellung der beteiligten Initiativen
  \item Mate, Wostok und Kekse
  \end{itemize}
\end{frame}

\section{Veranstaltungshinweise}
\begin{frame}
  \frametitle{Linux Helpdesk}
  \begin{itemize}
  \item LUG
  \item fsfw
  \end{itemize}
\end{frame}

%% \begin{frame}
%%   \frametitle{Interesse mitzumachen?}
%%   \begin{description}
%%   \item[Web] \url{https://fsfw-dresden.de}\\
%%     \begin{center}
%%       \hspace*{-6em}
%%       \includegraphics[width=4cm]{img-src/website-qr.pdf}
%%     \end{center}
%%   \item[Mailingliste] \url{http://lists.fsfw-dresden.de/mailman/listinfo/discuss} (auch oben verlinkt)
%%   \item[Treffen] SLUB (Hauptgebäude), Raum -2.115, Donnerstags in ungeraden Wochen, 18:30 Uhr
%%   \end{description}
%% \end{frame}

\end{document}
